\begin{abstract}
Windowed classification reduces sparse-host waiting under sustained interference
without sacrificing descriptor output service. In a unified register-transfer-level
(RTL) study of a single-output Compute Express Link (CXL) switch slice, the
FAI-inspired Main policy reduces sparse-host generation-to-fire P99 relative to
first-come, first-served (FCFS) by a mean paired 45.03\% for single arrivals and 38.64\% for
eight-request bursts. After two natural classification windows, both profiles
improve across all five matched seeds; all 15 production RTL runs pass.
All 30 measured intervals sustain the saturated
descriptor-level output capacity of 1.000000 transfer/cycle, with the heavy
host retaining approximately 97.7\% of service. Round robin with a ten-transfer
quantum (RR10) achieves a 3.6-cycle advantage in mean burst P99, while Main retains
a one-cycle advantage for single arrivals. Main's pre-layout area premium is
4.97--5.89\% over FCFS and 38.22--39.27\% over RR10. First-in, first-out
(FIFO) subtrees occupy 92.40--98.63\% of mapped standard-cell area; all nine
primary mappings close timing at 25\,ns. Independent audits and 14/14 scoped
formal RTL-to-netlist equivalence checks support the implementation comparison.
Classification shields sparse hosts relative to FCFS at a modest area
premium; per-source rotation yields lower measured burst P99 than Main.
Queue storage and ordering metadata dominate implementation cost.
\end{abstract}

\section{Introduction}
Arbitration can substantially reduce sparse-host waiting even when aggregate
output service is already saturated. Compute Express Link (CXL) enables memory
expansion and pooling through a coherent load/store interconnect~\cite{cxl_intro},
bringing heterogeneous host demands into competition at shared switch egress
ports. A continuously active host can saturate output capacity while occasional
requests wait behind its admitted work. Because work-conserving policies sustain
identical aggregate throughput under heavy load, aggregate bandwidth leaves the
central design trade-off unaddressed: which requests receive earlier service,
and what hardware cost buys that protection?

Our production register-transfer-level (RTL) experiments directly quantify this
exchange. Across 15 production runs, we evaluate first-come, first-served (FCFS),
round robin with a ten-transfer quantum (RR10), and FAI-inspired Main across
five matched seeds. After two natural classification windows, Main reduces
sparse-host P99 relative to FCFS by a mean paired 45.03\% for single arrivals and
38.64\% for bursts, improving in every seed under both profiles. All 30 measured
intervals sustain 1.000000 descriptor transfer per cycle, and the heavy host
retains approximately 97.7\% of service. Arbitration thus delivers earlier
sparse service at an unchanged total output rate.

Comparing Main against RR10 exposes the structural cost of aggregating sparse
sources into a single class. Main holds a one-cycle advantage in mean
single-arrival P99 (9.09\% relative to RR10), whereas RR10 achieves a 3.6-cycle
advantage in mean burst P99. Main's burst P99 is higher in four seeds and tied in
one. While RR10 grants an independent quantum to each source, Main's sparse
sources share a single ten-transfer High visit. Concurrent bursts can therefore
exhaust this shared class quota before Low service resumes. This mechanism is
consistent with the observed tail reversal, showing that the choice of service
domain remains consequential even at identical aggregate throughput.

We relate these service outcomes to mapped standard-cell area, timing, and
cold-start power within a common SystemVerilog switch. The synthesized boundary
integrates depth-16 per-source first-in, first-out (FIFO) queues, a 105-bit
descriptor datapath, one output, and a shared hold controller. Only policy
arbitration logic and necessary queue metadata vary. Across $N=4,8,16$ sources,
all nine primary mappings close timing at 25\,ns. Main's area premium is
4.97--5.89\% over FCFS and 38.22--39.27\% over RR10. Because FIFO subtrees
account for 92.40--98.63\% of total mapped area, accounting for storage and
ordering metadata is central to an accurate policy-cost comparison.

This study establishes a measured latency--area trade-off: two-window
classification protects sparse hosts relative to FCFS, while per-source
rotation remains preferable for the tested burst tails. It also shows why
selector-only area omits the dominant storage cost. Claims concern descriptor
arbitration and implementation cost at the evaluated switch boundary~\cite{artifact}.

\section{Related Work}
\paragraph{CXL sharing and interference.}
Pond investigates memory pooling for cloud platforms, TPP examines page placement
across local and CXL tiers, and DirectCXL demonstrates pooled memory
access~\cite{pond,tpp,directcxl}. These systems demonstrate the value of pooled
capacity, focusing on operating-system allocation, migration, and host-side
integration. Complementary studies examine switch architecture: A Case Against
CXL Memory Pooling and My CXL Pool analyze the economic and latency implications
of intermediate switches~\cite{case_pool,my_cxl_pool}. Furthermore,
CXL-INTERPLAY, contention-aware DRAM caching, and high-performance computing
(HPC) pooling studies characterize interference under shared
resources~\cite{cxl_interplay,contention_cache,hpc_pool}, while CXL-DMSim
connects full-system simulation with device measurements~\cite{cxldmsim}.
Our study isolates the descriptor datapath and measures how arbitration redistributes queuing delay.

\paragraph{Service history and grouping.}
Fair Queuing Memory Systems and stall-time fair scheduling distinguish service
allocation from application-level progress~\cite{fqmem,stfm}. ATLAS uses attained
service, and TCM groups threads by memory-access behavior~\cite{atlas,tcm}.
BLISS reduces scheduling complexity through two-group blacklisting and evaluates
Verilog implementations synthesized with Design Compiler and a commercial
32\,nm standard-cell library~\cite{bliss}. Its combination of workload behavior
and synthesized hardware cost serves as a methodological precedent for this report.
The grouping mechanisms differ: BLISS uses blacklisting, whereas Main uses
windowed successful-transfer counts and hysteresis that retains the previous
effective class.
Consistent with this literature, our study evaluates descriptor waiting time
and service allocation at the switch boundary rather than inferring
application-level IPC or slowdown.

\paragraph{CXL arbitration and hardware costs.}
FAI-CXL combines fairness-aware scheduling with cacheline-level
interleaving~\cite{fai}. We use its adaptive classification as the foundation
for an explicit single-output 10/10 reconstruction. The approved Main and a
supplementary Algorithm-2 interpretation feature distinct transition semantics,
enabling an empirical check on the implementation cost of that design choice.
Prior work in memory scheduling, crossbar arbitration, weighted-fair-queuing
sorting circuits, and network-on-chip (NoC) arbiters analyzes priority algorithms
and circuit complexity~\cite{rixner,gupta,wfq_sort,arbiter_ppa}. However, their
protocols, interfaces, and buffering models differ. In the literature reviewed
here, we found no directly comparable complete area, timing, and power baseline
at the same CXL switch boundary. Our common-boundary comparison provides this
baseline.

\section{Design and Service Mechanism}
\subsection{A Common Transfer Contract}
Figure~\ref{fig:architecture} illustrates the shared synthesized switch boundary. Each source
interfaces to a dedicated depth-16 FIFO. The 105-bit descriptor contains a 32-bit tag, 32-bit address,
one-bit opcode, eight-bit length, and 32-bit payload. All fields pass through the switch intact;
each descriptor consumes one atomic output service opportunity. FCFS and Main
also retain a 32-bit arrival field per FIFO entry. RR10 omits this ordering state
because its arbitration requires only a round-robin pointer and quantum counter. The resulting storage difference
is an integral component of each policy's actual implementation cost.

\begin{figure*}[t]
\centering
\resizebox{0.96\textwidth}{!}{%
\begin{tikzpicture}[font=\footnotesize,>=Latex,
  block/.style={draw=figgray,line width=.45pt,rounded corners=1pt,fill=white,align=center,inner sep=1.6pt},
  data/.style={-{Latex[length=1.2mm,width=.8mm]},draw=figgray,line width=.5pt},
  feedback/.style={-{Latex[length=1.2mm,width=.8mm]},draw=figred,line width=.65pt},
  note/.style={font=\scriptsize,text=figgray,align=center}]
  % Synthesized boundary
  \draw[dashed,figgray,line width=.45pt,rounded corners=2pt] (-.25,1.22) rectangle (17.10,6.70);
  \node[anchor=north west,fill=white,inner sep=1pt,font=\footnotesize\bfseries]
    at (-.18,6.62) {Synthesized RTL boundary};

  % Ingress lanes and per-source FIFOs
  \node[block,minimum width=1.55cm,minimum height=.72cm] (s0) at (1.00,5.95)
    {src 0\\[-1pt]R/V, 105b};
  \node[note] at (1.00,5.10) {$\vdots$};
  \node[block,minimum width=1.55cm,minimum height=.72cm] (sn) at (1.00,4.25)
    {src $N{-}1$\\[-1pt]R/V, 105b};
  \node[block,minimum width=1.70cm,minimum height=.72cm] (f0) at (3.20,5.95)
    {FIFO$_0$\\[-1pt]depth 16};
  \node[note] at (3.20,5.10) {$\vdots$};
  \node[block,minimum width=1.70cm,minimum height=.72cm] (fn) at (3.20,4.25)
    {FIFO$_{N-1}$\\[-1pt]depth 16};
  \draw[data] (s0) -- (f0);
  \draw[data] (sn) -- (fn);

  % FIFO-head candidate bus
  \draw[data] (f0.east) -- (4.55,5.95) -- (4.55,5.28);
  \draw[data] (fn.east) -- (4.55,4.25) -- (4.55,5.28);
  \draw[data] (4.55,5.28) -- (5.18,5.28);
  \node[note,fill=white,inner sep=1pt,anchor=south] at (4.55,5.68) {heads};

  % One elaboration-time policy slot plus its retained context
  \draw[draw=figgray,line width=.55pt,rounded corners=1.5pt,fill=figblue!4]
    (5.18,3.35) rectangle (8.85,6.45);
  \node[font=\footnotesize\bfseries] at (7.015,6.20) {Policy slot};
  \node[note] at (7.015,5.93) {one selected at elaboration};
  \node[block,draw=figblue!75,fill=figblue!9,minimum width=1.25cm,minimum height=.44cm,inner sep=1pt]
    at (6.27,5.52) {FCFS};
  \node[block,draw=figblue!75,fill=figblue!9,minimum width=1.25cm,minimum height=.44cm,inner sep=1pt]
    at (7.75,5.52) {RR};
  \node[block,draw=figred!75,fill=figred!7,minimum width=1.25cm,minimum height=.44cm,inner sep=1pt]
    at (6.27,4.98) {Main};
  \node[block,draw=figred!75,fill=figred!7,minimum width=1.25cm,minimum height=.44cm,inner sep=1pt]
    at (7.75,4.98) {Alg-2};
  \draw[figgray!55,dashed,line width=.4pt] (5.35,4.58) -- (8.68,4.58);
  \node[block,draw=figgray!55,fill=white,minimum width=2.75cm,minimum height=.86cm,inner sep=1pt]
    at (7.015,4.05) {candidate source\\+ policy context};

  % Common datapath and output
  \draw[data] (8.85,5.28) -- (9.23,5.28);
  \draw[draw=figgray,line width=.5pt,fill=white]
    (9.23,4.75) -- (9.23,5.81) -- (10.68,5.55) -- (10.68,5.01) -- cycle;
  \node[font=\footnotesize\bfseries] at (9.90,5.28) {MUX};
  \node[block,minimum width=1.85cm,minimum height=.74cm] (hold) at (11.75,5.28)
    {source/context\\[-1pt]hold};
  \draw[data] (10.68,5.28) -- (10.80,5.28);
  \draw[data] (hold.east) -- (16.78,5.28);
  \node[note,above] at (14.90,5.38) {105b R/V};
  \node[note,anchor=east] at (16.68,5.10) {to CXL};

  % Fire controls FIFO pop, quota use, and held context.
  \draw[feedback] (16.78,5.28) -- (16.78,2.62) -- (3.20,2.62);
  \draw[feedback] (3.20,2.62) -- (3.20,3.88);
  \draw[feedback] (7.02,2.62) -- (7.02,3.33);
  \node[note,fill=white,inner sep=2pt,anchor=south] at (11.35,2.77)
    {\code{fire} = valid $\land$ ready};
  \node[note,fill=white,inner sep=2pt,anchor=south] at (5.85,2.77)
    {pop + quota};

  % Descriptor layout is called out separately; it is not to scale.
  \draw[draw=figgray!65,line width=.4pt,rounded corners=1pt,fill=white]
    (0.18,1.35) rectangle (8.95,2.23);
  \node[anchor=south west,font=\scriptsize,fill=white,inner sep=1pt]
    at (.30,2.26) {FIFO entry (not to scale)};
  \foreach \x/\w in {.33/1.26,1.59/1.26,2.85/.78,3.63/1.00,4.63/1.43,6.06/1.58}
    \draw[figgray!55,line width=.35pt] (\x,1.43) rectangle ++(\w,.55);
  \node[note] at (.96,1.705) {Tag\\[-1pt][104:73]};
  \node[note] at (2.22,1.705) {Addr.\\[-1pt][72:41]};
  \node[note] at (3.24,1.705) {Opcode\\[-1pt][40]};
  \node[note] at (4.13,1.705) {Length\\[-1pt][39:32]};
  \node[note] at (5.345,1.705) {Payload\\[-1pt][31:0]};
  \draw[dashed,figred!75,line width=.45pt] (6.06,1.43) rectangle (7.64,1.98);
  \node[note] at (6.85,1.705) {Arrival\\[-1pt][31:0]$^{*}$};
  \node[note,anchor=west] at (9.18,1.70)
    {$^{*}$Arrival metadata is present only for\\FCFS, Main, and Algorithm-2.};

  % Non-synthesized testbench
  \draw[dashed,figgray,line width=.45pt,rounded corners=2pt,fill=figgray!2]
    (-.25,.08) rectangle (17.10,1.12);
  \node[anchor=north west,fill=figgray!2,inner sep=1pt,font=\footnotesize\bfseries]
    at (-.18,1.08) {Testbench only (not synthesized)};
  \node[block,minimum width=2.55cm,minimum height=.50cm] at (1.45,.47) {Traffic producers};
  \node[block,minimum width=2.65cm,minimum height=.50cm] at (5.00,.47) {Memory-ready trace};
  \node[block,minimum width=2.55cm,minimum height=.50cm] at (8.80,.47) {Python oracle};
  \node[block,minimum width=2.75cm,minimum height=.50cm] at (12.75,.47) {Response scoreboard};
\end{tikzpicture}
}
\caption{Controlled synthesis boundary. The selected policy proposes a source independently of downstream readiness. The orange path denotes \code{fire}; it pops the selected FIFO, consumes service quota, and preserves the source/context hold through a stall. The testbench drives and observes DUT pins; its Python oracle is an independent reference model, and the scoreboard checks observed responses. These blocks are not synthesized.}
\label{fig:architecture}
\end{figure*}


A push is accepted whenever beginning-of-cycle FIFO occupancy is below 16.
The arbiter selects an eligible source independently of sink readiness, and
\begin{equation}
 f(t)=v_{\rm out}(t)\land r_{\rm out}(t)
 \label{eq:fire}
\end{equation}
defines successful service. Only $f(t)$ pops a FIFO and decrements active quota. The common
controller holds the selected source, descriptor, and policy context across a
downstream stall. Per-source ordering, admission rules, and output holding are
identical across policies.

\subsection{Capacity Is Shared; Waiting Is Scheduled}
For a single output serving equal-cost atomic descriptors, the measured service rate is
\begin{equation}
 R=\frac{1}{T}\sum_{t=0}^{T-1}f(t)\leq 1.
 \label{eq:capacity}
\end{equation}
Since $f(t)$ is binary, $R=1$ requires a successful transfer in every measured
cycle. A work-conserving arbiter sustains this rate whenever each cycle
presents eligible FIFO work and a ready receiver ($r_{\rm out}(t)=1$). External
offered demand alone is insufficient: producer backlog does not constitute eligible work,
and a newly accepted descriptor cannot bypass unread FIFO entries. Under continuous occupancy
and downstream readiness, the output's descriptor service capacity is fixed at 1.000000 transfer
per cycle. Arbitration can still change which descriptor fires, its waiting time, and
each source's service share. Saturated throughput therefore leaves waiting
tails policy-dependent; maximum throughput does not establish equivalent service quality.

FCFS consistently schedules the oldest admitted queue head, with source-ID tie breaking. Sparse
requests therefore compete directly with older requests already admitted from the heavy
source. Depth-16 FIFOs constrain this admitted competition, while unadmitted
heavy requests accumulate in the external producer queue. RR10 instead rotates across
nonempty sources, granting up to ten successful transfers per visit and skipping
empty sources. It organizes competition by source identity rather than arrival
age. Downstream sink stalls consume no quantum under either quota-based policy.

\subsection{Classification Changes the Service Domain}
Main tracks successful transfers $B[h]$ for each source across a $W$-cycle window,
computing total delivered service $S=\sum_h B[h]$. Its classification thresholds evaluate
\begin{align}
 \mathrm{UP}:&\quad 10NB[h]<3S,\nonumber\\
 \mathrm{DOWN}:&\quad NB[h]>S,\label{eq:classes}\\
 \mathrm{KEEP}:&\quad \text{otherwise, including equality}.\nonumber
\end{align}
Production uses $W=5{,}000{,}000$, 23-bit counters, and widened unsigned 32-bit
comparisons. The window-closing transfer is included before counters clear. Every
source initializes in effective High. Two supporting windows commit a transition;
a pending state retains its previous effective class. Observing KEEP, opposite evidence,
or an idle window cancels the pending transition. Classification thus reflects
sustained delivered service over time rather than transient arrival spikes.

Main alternates visits between High and Low classes, each with a quota of up to ten fires,
and selects the oldest head within the active class. Under heavy/sparse traffic,
two windows move the heavy source to Low while sparse sources remain High.
Older heavy entries then compete in a separate service domain from sparse
entries. A sparse request still waits for active visit context, but its
service no longer depends on outranking every heavy head by arrival age.

An empty High visit immediately relinquishes unused quota. That credit is not
banked, and Low service proceeds without an artificial idle cycle. The 10/10 schedule
therefore divides output cycles only when both classes need service; it never reserves
capacity for absent sparse work. This mechanism explains how Main preserves
full 1.000000 transfer/cycle output service while the heavy host uses approximately 97.7\% of it.

However, class aggregation creates competition within the protected class. Sparse sources
share one ten-fire High visit, whereas RR10 assigns ten fires to each source.
Overlapping sparse bursts can exhaust the shared High visit and wait through
an intervening Low visit. This structure can lengthen sparse waits despite full output
service, motivating separate comparisons of Main against RR10 and FCFS.

\section{Implementation and Experimental Method}
\subsection{Correctness and Comparable Hardware}
A single policy-selectable SystemVerilog top-level module shares FIFO, multiplexer, and hold logic.
Independent golden Python queue models and cocotb testbenches verify cycle-accurate output pin activity,
descriptors, source order, conservation, response tags, buffer occupancies, and held-output stability.
Directed tests exercise timestamp wrap, full-plus-pop refusal, bypass prevention, quota
boundaries, class thresholds, and reset assertion during output holding. The accepted
primary regression achieves 180/180 passing runs across three policies, three source
counts, and 20 seeds. Main additionally passes 30 stress runs and three
20-million-cycle production runs across $N=4,8,16$.

Synopsys Formality W-2024.09 formally proves 14/14 scoped RTL/netlist pairs: nine primary
mappings, two supplementary mappings, and three current-Main bridges. The verified pairs
have zero failing, aborted, unverified, unmatched, or active-blackbox points.
An intentional output-valid inversion fails its expected compare point,
validating our negative control. Six FCFS/RR pairs retain one explicitly
uncompared propagated-constant register. The supplementary Algorithm-2 study
passes ten stress cases and a production run; at $N=4$ its complete-area
difference from Main is 0.028626\%. We retain it as transition-semantic
sensitivity evidence and focus the primary comparison on FCFS, RR10, and Main.

The nine primary area/timing mappings use Design Compiler and PrimeTime
W-2024.09 with SAED32 RVT standard cells at the TT corner (0.78\,V, $25\,{}^\circ$C).
The common clock is 25\,ns, setup/hold uncertainty is 0.10/0.05\,ns, input/output delays are
0/2\,ns, and the nominal output load is 2.656828\,fF. A common wire-load
estimate and identical constraints define the pre-layout comparison. Final
checks restore the original output capacitance limits after optimization.
Area is reported in library units. The accepted cold-power matrix uses
zero-delay VCS mapped-netlist replay and switching activity interchange format
(SAIF) data in PrimeTime/PrimePower
at the same corner and clock, covering three seeds per mapping.

\subsection{Production HOG-SPARSE Experiment}
The accepted HOG-SPARSE benchmark comprises $3\times5=15$ actual RTL runs at
$N=4$, depth 16, with seeds 509--513 and a continuously ready sink. Source zero
generates one descriptor per active cycle. Sources one through three exhibit
sparse demand, driven by deterministic, policy-independent pseudorandom streams.
Each run executes ten million real training cycles, covering two complete
production windows. Main's heavy source transitions through High, pending High, then Low;
the sparse sources remain effective High. The measured class map is
Low/High/High/High.

Profile A continues independent Bernoulli arrivals with probability $1/128$
per sparse source per cycle for 100,000 measured cycles. Profile B generates
one eight-descriptor burst per sparse source per 1,024-cycle block, at a seeded
uniform offset, for 102,400 cycles. Its 100 blocks generate exactly 2,400 sparse
requests. Both profiles maintain expected sparse demand of $3/128$ requests/cycle,
allowing temporal concentration to be examined at the same expected offered
rate; Profile A retains its realized finite counts. The heavy source supplies persistent
competition throughout training, both profiles, and completion guards.

Each profile is followed by a 1,024-cycle guard that stops new sparse arrivals
while maintaining heavy generation. Every sparse cohort thus completes with heavy
pressure still active. Final drain then flushes all producer, FIFO, and
response work. The accepted execution accounts for 156,657,531 nonreset RTL
cycles, including 150 million training cycles and 3,036,000 measured cycles.
All 15 production runs pass, and an independent raw audit passes. The audit
checks exact histograms, source pooling, direct/stage identities, conservation,
class transitions, completion, and actual run exits. Five external trace groups
match across policies; Main and FCFS also produce identical cold pin digests over
each seed's entire ten-million-cycle training interval.

\subsection{Endpoints and Paired Summaries}
For request $i$, let $g_i,a_i,d_i$ denote its generation, acceptance, and fire cycles.
The primary endpoint is direct sparse latency
\begin{equation}
 L_i=d_i-g_i=(a_i-g_i)+(d_i-a_i).
 \label{eq:latency}
\end{equation}
Each run pools the complete sparse generation cohort of one profile and extracts
nearest-rank P99 directly from its request histogram. Per-source results remain
available for checking distribution within the class. Figures report arithmetic
means of five run means or five run P99s; whiskers denote observed seed ranges.
For baseline $b$ and matched seed $s$, the percentage reduction is
\begin{equation}
 \rho_{b,s}=100\frac{P99_{b,s}-P99_{{\rm Main},s}}{P99_{b,s}}.
 \label{eq:paired}
\end{equation}
We average these five paired percentages, retaining the baseline denominator
within each seed. This differs from a ratio of mean P99s; negative values
denote higher Main P99. The five seeds provide descriptive replications;
sources sharing an arbiter are not independent replications. Seed ranges
are neither confidence intervals nor worst-case bounds.

\section{Evaluation}
\subsection{Sparse Protection at Full Output Capacity}
Main protects sparse single arrivals while the descriptor output remains fully
saturated (Figure~\ref{fig:hog-sparse}). In Profile A, FCFS exhibits a mean direct
latency of \HSMeanAFCFS\ cycles and a mean run P99 of \HSPtailAFCFS\ cycles;
Main reduces these to \HSMeanAMain\ and \HSPtailAMain\ cycles.
The P99 reduction averages 8.2 cycles and \HSReductionAFCFS\% across matched
seeds, improving in five of five. RR10 achieves \HSMeanARR\ cycles mean and
\HSPtailARR\ cycles mean run P99. Main's one-cycle P99 advantage over RR10
is \HSReductionARR\%, likewise in every seed.

\begin{figure*}[t]
\centering
\begin{tikzpicture}
\begin{groupplot}[
  group style={group size=2 by 2,horizontal sep=5.5em,vertical sep=4.8em},
  width=.435\textwidth,height=3.65cm,
  title style={font=\small},tick label style={font=\footnotesize,text=figgray},
  label style={font=\footnotesize},
  axis line style={figgray!70},tick style={figgray!70},
  axis x line=bottom,axis y line=left,
  ymajorgrids,grid style={figgray!18,line width=.25pt},
  xmin=.55,xmax=2.45,xtick={1,2},xticklabels={A: single arrivals,B: eight-request bursts},
  xlabel={Sparse arrival profile},
  /pgfplots/error bars/y dir=both,/pgfplots/error bars/y explicit,
  /pgfplots/error bars/error bar style={line width=.6pt},
  /pgfplots/error bars/error mark options={mark size=2pt}]
\nextgroupplot[title={(a) Sparse mean latency},ylabel={Generation-to-fire (cycles)},
  ymin=0,ymax=28,ytick={0,5,10,15,20,25},ybar,bar width=10pt]
\addplot[fill=figblue!65,draw=figblue,bar shift=-11pt] table[x=x,y=y,y error minus=lo,y error plus=hi]{\HSMeanFCFS};
\addplot[fill=figgray!25,draw=figgray,dashed,bar shift=0pt] table[x=x,y=y,y error minus=lo,y error plus=hi]{\HSMeanRR};
\addplot[fill=figred!65,draw=figred,bar shift=11pt] table[x=x,y=y,y error minus=lo,y error plus=hi]{\HSMeanMain};
\nextgroupplot[title={(b) Sparse P99 latency},ylabel={Generation-to-fire P99 (cycles)},
  ymin=0,ymax=40,ytick={0,10,20,30,40},ybar,bar width=10pt]
\addplot[fill=figblue!65,draw=figblue,bar shift=-11pt] table[x=x,y=y,y error minus=lo,y error plus=hi]{\HSPtailFCFS};
\addplot[fill=figgray!25,draw=figgray,dashed,bar shift=0pt] table[x=x,y=y,y error minus=lo,y error plus=hi]{\HSPtailRR};
\addplot[fill=figred!65,draw=figred,bar shift=11pt] table[x=x,y=y,y error minus=lo,y error plus=hi]{\HSPtailMain};
\nextgroupplot[title={(c) Main's paired P99 change},ylabel={P99 reduction (\%)},
  xmin=.5,xmax=4.5,xtick={1,2,3,4},xticklabels={A/FCFS,A/RR10,B/FCFS,B/RR10},
  xlabel={Profile / comparison baseline},ymin=-50,ymax=60,ytick={-40,-20,0,20,40,60}]
\addplot[figgray,densely dashed,no marks,forget plot] coordinates {(.5,0)(4.5,0)};
\addplot[only marks,figgray,mark=*,mark size=1.15pt,forget plot] coordinates {\HSPairedSeeds};
\addplot[only marks,figred,mark=diamond*,mark size=3pt] table[x=x,y=y,y error minus=lo,y error plus=hi]{\HSPaired};
\nextgroupplot[title={(d) Heavy service at full capacity},ylabel={Transfers / measured cycle},
  ymin=0,ymax=1.12,ytick={0,.25,.5,.75,1},ybar,bar width=10pt]
\addplot[fill=figblue!65,draw=figblue,bar shift=-11pt] table[x=x,y=y,y error minus=lo,y error plus=hi]{\HSHeavyFCFS};
\addplot[fill=figgray!25,draw=figgray,dashed,bar shift=0pt] table[x=x,y=y,y error minus=lo,y error plus=hi]{\HSHeavyRR};
\addplot[fill=figred!65,draw=figred,bar shift=11pt] table[x=x,y=y,y error minus=lo,y error plus=hi]{\HSHeavyMain};
\addplot[black,densely dashed,line width=.7pt,no marks,sharp plot,forget plot] coordinates {(.55,1)(2.45,1)};
\end{groupplot}
\node[anchor=south,font=\footnotesize,yshift=31pt] at
  ($(group c1r1.north west)!.5!(group c2r1.north east)$)
  {\tikz\draw[fill=figblue!65,draw=figblue] (0,0) rectangle (.22,.12);\ FCFS\qquad
   \tikz\draw[fill=figgray!25,draw=figgray,dashed] (0,0) rectangle (.22,.12);\ RR10\qquad
   \tikz\draw[fill=figred!65,draw=figred] (0,0) rectangle (.22,.12);\ Main};
\end{tikzpicture}
\caption{\textbf{Sparse-host protection with unchanged output capacity.}
Fifteen production RTL runs, $N=4$, seeds 509--513, two real $W=5$M training windows per run.
(a,b) Bars are equal-weight means of five run means and five run P99s, respectively;
whiskers span the observed seed minimum and maximum, not confidence intervals.
(c) Diamonds are mean paired percentage reductions; dots are the five matched seeds
(horizontal jitter is for visibility). Positive values favor Main.
(d) Bars show heavy-host service; the dashed line is total output service,
exactly 1.000000 transfer/cycle in every measured interval. Heavy share remains
97.66--97.67\%; its seed ranges are smaller than the plotting resolution.
Profile A lasts 100,000 cycles; B lasts 102,400 cycles.}
\label{fig:hog-sparse}
\end{figure*}


Profile B preserves Main's advantage over FCFS under burst injection. FCFS
reaches \HSMeanBFCFS\ cycles mean and \HSPtailBFCFS\ cycles P99; Main reaches
\HSMeanBMain\ and \HSPtailBMain\ cycles. The paired reduction is 13.0 cycles
or \HSReductionBFCFS\%, with all five seeds improving. Across the three sparse
sources and five seeds, every one of the 15 source/seed P99 comparisons improves
over FCFS in each profile. The pooled improvement thus extends to each individual sparse
source in the tested runs, rather than masking a source-level regression.

Stage-latency decomposition locates the benefit after FIFO admission. In Profile A,
the generation-to-acceptance mean is zero under every policy, proving that earlier output
service accounts entirely for Main's improvement. In Profile B, all policies incur the same
3.5-cycle mean injection delay: eight descriptors generated simultaneously enter a
source FIFO at one per cycle. Main's acceptance-to-fire mean is 5.8243 cycles,
versus FCFS's 19.7989 cycles. The reduction therefore survives both burst injection
and its common serialization cost. Direct P99 is measured from each request's
full generation-to-fire wait; stage percentiles are never added.

The cold control supports the class-separation mechanism. Before the second
classification window, Main and FCFS produce matching per-seed sparse
histograms, a five-run mean of 16.3842 cycles, and P99 of 18 cycles; their full
training pin digests also match. Main's protection follows its two-window
transition. Both profiles finish before a third production window.

\subsection{The RR10 Burst Trade-off}
RR10 demonstrates a 3.6-cycle advantage in mean burst P99: \HSPtailBRR\ cycles versus
Main's \HSPtailBMain\ cycles. Its mean direct latency is \HSMeanBRR\ cycles, close
to Main's \HSMeanBMain\ cycles, showing that a negligible mean difference can conceal a larger tail
divergence. Using RR10 as the baseline denominator in each paired comparison, Main's
burst P99 is 20.78\% higher on average; Equation~\ref{eq:paired} therefore
returns $-20.78\%$. Main exhibits higher P99 in four seeds and ties in one, with
run P99s spanning 17--24 cycles versus RR10's tighter 17--18 cycle range.

The service domains defined in Section~3 explain why an improvement over FCFS does
not imply an improvement over RR10. Both Main and RR10 eliminate global
arrival-age competition against the heavy host. Overlapping sparse bursts,
however, share a single High visit in Main, whereas RR10 retains an independent quantum per source.
Class aggregation is consistent with the measured tail difference, although
no separate causal ablation isolates it. For this source mix, Main retains
a one-cycle single-arrival P99 advantage, whereas RR10 delivers lower measured
burst P99 while occupying less silicon area.

\subsection{Heavy Service and Conservation}
All 30 measured policy/seed/profile intervals sustain exactly 1.000000
descriptor transfer/cycle with zero avoidable eligible-ready bubbles. In Profile A,
heavy transfers/cycle average 0.976728 for FCFS, 0.976726 for RR10, and 0.976724
for Main. In Profile B, the respective rates are 0.97657617, 0.97656445, and
0.97656445. Because total service is 1.000000, these rates represent the heavy
host's service fractions: approximately 97.66--97.67\%.

Main's mean heavy-rate difference from FCFS in Profile A is only $-0.000004$
transfer/cycle. Empty-High relinquishment lets the heavy Low source consume
spare opportunities, ensuring that sparse protection does not impose an artificial 50\% heavy-host
cap. The small interval differences include requests crossing phase boundaries;
complete-cohort latency tracking and final conservation checks account for all such requests.

Total offered load exceeds descriptor output capacity: the heavy host alone
offers one request per cycle, with sparse work added. Its producer backlog
grows under every policy; at the end of Profile B,
mean pending counts are 239,330.4 for FCFS and 239,331.6 for RR10 and Main.
These counts describe unadmitted producer work, separate from each depth-16
FIFO. Final drain accounts for every generated heavy request, while the sparse
completion guards establish protection during continuing interference. These
overloaded intervals do not define a stationary heavy-host delay distribution.

\subsection{Policy-to-Hardware Cost}
Figure~\ref{fig:hardware-cost} juxtaposes the service results with the accepted
hardware comparison. All nine primary mappings meet setup, hold, and original
electrical limits at the common 25\,ns target. Main adds 5.89\%, 5.23\%, and
4.97\% total area over FCFS at $N=4,8,16$, respectively. Its premiums over
RR10 are 39.27\%, 38.60\%, and 38.22\%. At the four-source boundary used
for HOG-SPARSE, Main's area is 110,089.846 library units, compared with
103,970.058 for FCFS and 79,047.172 for RR10. The measured sparse-host
protection over FCFS therefore accompanies a 5.89\% complete-area premium
at the same source count.

\begin{figure}[!t]
\centering
\begin{tikzpicture}
\begin{groupplot}[
  group style={group size=1 by 4,vertical sep=4.5em},
  width=.86\columnwidth,height=2.3cm,
  title style={font=\scriptsize},tick label style={font=\tiny,text=figgray},
  label style={font=\scriptsize},axis line style={figgray!70},tick style={figgray!70},
  grid=major,grid style={figgray!15,line width=.25pt},
  axis x line=bottom,axis y line=left,
  xmin=2.5,xmax=17.5,xtick={4,8,16},xlabel={Sources $N$}]
\nextgroupplot[title={(a) Complete-switch area},ylabel={Area ($10^3$ library units)},
  ymin=0,ymax=480,ytick={0,100,200,300,400}]
\addplot[figblue,line width=.8pt,mark=square*,mark size=1.8pt] coordinates {\AreaFCFS};
\addplot[figgray,line width=.8pt,dashed,mark=triangle*,mark options={fill=white},mark size=2pt] coordinates {\AreaRR};
\addplot[figred,line width=.8pt,mark=*,mark size=1.8pt] coordinates {\AreaMain};
\nextgroupplot[title={(b) Timing margin},ylabel={Setup slack (ns)},
  ymin=0,ymax=23,ytick={0,5,10,15,20}]
\addplot[figblue,line width=.8pt,mark=square*,mark size=1.8pt] coordinates {\SetupFCFS};
\addplot[figgray,line width=.8pt,dashed,mark=triangle*,mark options={fill=white},mark size=2pt] coordinates {\SetupRR};
\addplot[figred,line width=.8pt,mark=*,mark size=1.8pt] coordinates {\SetupMain};
\nextgroupplot[title={(c) Cold-start power},ylabel={Power (mW)},
  ymin=0,ymax=16,ytick={0,4,8,12,16}]
\addplot[figblue,line width=.8pt,mark=square*,mark size=1.8pt] coordinates {\ColdMeanFCFS};
\addplot[figgray,line width=.8pt,dashed,mark=triangle*,mark options={fill=white},mark size=2pt] coordinates {\ColdMeanRR};
\addplot[figred,line width=.8pt,mark=*,mark size=1.8pt] coordinates {\ColdMeanMain};
\nextgroupplot[title={(d) Service-cost frontier},
  ylabel={P99 (cycles)},
  xmin=73,xmax=115,xtick={80,90,100,110},
  xlabel={Mapped area ($10^3$ units)},
  ymin=0,ymax=40,ytick={0,10,20,30,40}]
\addplot[figblue,line width=.7pt,dotted,forget plot] coordinates
  {(103.970058,18.2)(103.970058,33.8)};
\addplot[figgray,line width=.7pt,dotted,forget plot] coordinates
  {(79.047172,11.0)(79.047172,17.2)};
\addplot[figred,line width=.7pt,dotted,forget plot] coordinates
  {(110.089846,10.0)(110.089846,20.8)};
\addplot[only marks,figblue,mark=square*,mark size=2.1pt] coordinates {(103.970058,33.8)};
\addplot[only marks,figblue,mark=square,mark options={fill=white},mark size=2.1pt]
  coordinates {(103.970058,18.2)};
\addplot[only marks,figgray,mark=triangle*,mark size=2.3pt] coordinates {(79.047172,17.2)};
\addplot[only marks,figgray,mark=triangle,mark options={fill=white},mark size=2.3pt]
  coordinates {(79.047172,11.0)};
\addplot[only marks,figred,mark=*,mark size=2.1pt] coordinates {(110.089846,20.8)};
\addplot[only marks,figred,mark=o,mark options={fill=white},mark size=2.1pt]
  coordinates {(110.089846,10.0)};
\node[font=\tiny,anchor=north,text=figblue] at (axis cs:103.970058,17.2) {FCFS};
\node[font=\tiny,anchor=north,text=figgray] at (axis cs:79.047172,10.0) {RR10};
\node[font=\tiny,anchor=north,text=figred] at (axis cs:110.089846,9.0) {Main};
\end{groupplot}
\node[anchor=south,font=\scriptsize,yshift=15pt] at
  ($(group c1r1.north west)!.5!(group c1r1.north east)$)
  {\tikz\fill[figblue] (0,0) rectangle (.09,.09);\ FCFS\qquad
   \tikz\draw[figgray] (0,0)--(.045,.08)--(.09,0)--cycle;\ RR10\qquad
   \tikz\fill[figred] (0,0) circle (.045);\ Main\qquad
   open: A; filled: B};
\end{tikzpicture}
\caption{\textbf{The hardware price of service differentiation.}
Accepted SAED32 RVT mappings at TT, 0.78\,V, $25\,{}^\circ$C and a common
25\,ns clock. (a,b) One mapped observation per policy and $N$.
(c) Means of three cold-start seed estimates.
(d) $N=4$ P99 versus complete mapped area; open/filled markers denote
single arrivals (A) and eight-request bursts (B). Lower-left is better.
Power uses a separate zero-delay mapped-netlist activity experiment.}
\label{fig:hardware-cost}
\end{figure}


FIFO subtrees account for 92.40--98.63\% of complete mapped area, including
arrival metadata where required. At $N=16$, 256 FIFO entries hold 26,880
descriptor bits; FCFS and Main add 8,192 arrival bits. RR10 avoids this
per-entry ordering state, so its smaller footprint reflects the service rule
as well as simpler arbitration logic. Main further
adds window counters, class state, and quota context. Its policy subtree
grows from 7,143.988 library units at $N=4$ to 26,673.175 at $N=16$,
whereas FCFS grows from 1,131.957 to 6,038.461 and RR10 from 263.039
to 1,268.687. Comparing these policy subtrees alone would omit the FIFO
metadata cost. The complete-area result charges both, exposing storage as
the dominant budget even when selector complexity differs substantially.

Timing closure exhibits a policy-dependent margin.
FCFS setup slack decreases from 13.072\,ns at $N=4$ to 3.892\,ns at
$N=16$; Main decreases from 10.496 to 1.158\,ns. RR10 retains at least
18.694\,ns across all configurations. Main16's limiting timing path connects FIFO read-pointer state to
class-state update, linking head selection to the closing-window decision.
All reported designs use the common pipeline and constraints. The accepted
timing accounting retains 140 constant-disabled checks separately, so closure
refers to evaluated setup, hold, and electrical checks under this profile.

Cold-start power consistently orders RR10 $<$ FCFS $<$ Main for every
accepted seed. The three-seed means at $N=4$ are approximately 2.601,
3.398, and 3.500\,mW; at $N=16$ they are 10.3, 13.4, and 13.8\,mW. These
estimates use their own cold activity trace, separate from the production
HOG-SPARSE latency experiment.

\section{Conclusion}

Main reduces sparse-host P99 relative to FCFS by a mean paired 45.03\% for
single arrivals and 38.64\% for bursts after two natural classification
windows, improving in every matched seed. All 30 measured intervals sustain
1.000000 descriptor transfer/cycle; the heavy host retains 97.66--97.67\%
of service because empty-High relinquishment releases unused quota.
Arbitration therefore changes who waits even when descriptor output service
is saturated. RR10 has a 3.6-cycle advantage in mean burst P99, while Main
retains a one-cycle advantage for single arrivals. Main's area premium over
RR10 is 38.22--39.27\%. At $N=4$, Figure~\ref{fig:hardware-cost}(d) makes
the choice concrete: Main reduces single-arrival P99 by 8.2 cycles over FCFS
for 5.89\% more area, whereas RR10 is smaller than Main and has lower
measured burst P99.

Main's protection over FCFS costs 4.97--5.89\% total area across $N=4,8,16$.
All nine primary mappings close timing at 25\,ns, and 14/14 scoped
RTL-to-netlist equivalence checks support implementation fidelity. FIFO
subtrees occupy 92.40--98.63\% of mapped standard-cell area, so storage
dominates implementation cost while service order governs waiting tails.

\section{AI Use Acknowledgment}
Claude Code/RTL Agent Team and owner-authorized GPT-6.1 Sol supported implementation
and verification tooling; Codex coordinated evidence review and drafting under ARS-Codex guidance.
Gemini 4 Argon assisted literature analysis and claim review.
Gemini 3.8 Flash performed the final academic language polish.
Numerical findings derive strictly from archived tool results.
The authors assume full responsibility for this report and its sources.

